% SPDX-License-Identifier: CC-BY-SA-4.0
% Part: Syntactic analysis
% Section: Cocke-Younger-Kasami algorithm
% Challenge exercise: Greibach normal form

\begingroup

\begin{exercice}[Mini-défi -- Forme normale de Greibach]\label{exo:parsing/CYK/challenge}

  Une grammaire algébrique $G=\langle \Sigma, \Gamma, S, \rightarrow \rangle$ est en \emph{forme normale de Greibach} si toute règle est de la forme
  $A \rightarrow a X_1 X_2 \cdots X_k$, où $A \in \Gamma$, $a \in \Sigma$, $k\in \mathbb{N}$ et pour tout $i\in \{1, ..., k\}$, $X_i \in \Gamma$. 
  La règle $S\rightarrow \varepsilon$ est autorisée \emph{uniquement} si $\varepsilon\in \mathcal{L}(G)$.

  \begin{question}

  \item Pour chacune des grammaires suivantes, donnez une grammaire sous forme normale de Greibach générant le même langage. 
    \begin{itemize}

    \item $G_1 \eqdef \left\langle \{a,b\}, \{S,T\}, S,
      \left\{\begin{array}{l}
      S \rightarrow a S \mid T\\
      T \rightarrow bT \mid \varepsilon
      \end{array}\right\}\right\rangle$

    \item $G_2 \eqdef \left\langle \{a,b\}, \{S,A\}, S,
      \left\{\begin{array}{l}
      S \rightarrow A S \mid a\\
      A \rightarrow b
      \end{array}\right\}\right\rangle$

    \item $G_3 \eqdef \left\langle \{a,b\}, \{S\}, S,
      \left\{\begin{array}{rcl}
      S &\rightarrow& S b \mid a
      \end{array}\right\}\right\rangle$


    \end{itemize}
    
  \item Proposer un algorithme qui transforme toute grammaire en une grammaire
    équivalente en forme normale de Greibach.

  \item Appliquer la méthode à la grammaire suivante : 
    $\left\langle \{a,b,c\}, \{A,B,C\}, A,
    \left\{\begin{array}{rcl}
    A & \rightarrow & AB\mid a\\
    B & \rightarrow & BC\mid b\\
    C & \rightarrow & CA\mid c\\
    \end{array}\right\}\right\rangle$

  \end{question}

\end{exercice}


\begin{correction}

  \begin{question}

  \item $G_1$ engendre $a^\star b^\star$, $G_2$ engendre $b a^+$ et $G_3$ engendre $a b^+$. Ces langages sont reconnus par les grammaires suivantes sous forme normale de Greibach. 
    \begin{itemize}

    \item $G_1 \eqdef \left\langle \{a,b\}, \{S,T\}, S,
      \left\{\begin{array}{l}
      S \rightarrow a S \mid b T \mid \varepsilon\\
      T \rightarrow b T \mid b
      \end{array}\right\}\right\rangle$

    \item $G_2 \eqdef \left\langle \{a,b\}, \{S,A\}, S,
      \left\{\begin{array}{l}
      S \rightarrow b S \mid a
      \end{array}\right\}\right\rangle$

    \item $G_3 \eqdef \left\langle \{a,b\}, \{S\}, S,
      \left\{\begin{array}{rcl}
      S &\rightarrow& a T b \\
      T &\rightarrow& b T
      \end{array}\right\}\right\rangle$

    \end{itemize}
    
  \item Esquisse d'algorithme vers la FNG.
    \begin{enumerate}
    \item \emph{Prétraitements.} Éliminer les symboles inutiles, les $\varepsilon$-productions,
      les règles unitaires ($A\rightarrow B$ avec $A\neq B$),
      et introduire des règles $T_a\rightarrow a$ pour chaque terminal $a\in\Sigma$, 
      comme pour la mise sous forme normale de Chosky.
    \item On élimine chaque chaque non-terminal $A$ en début de membre gauche de règle, un par un : 
      \begin{enumerate}
      \item \emph{Élimination de la récursion gauche.}
        Les règles de la forme $A \rightarrow A\alpha \mid \beta$ sont remplacées par $A \rightarrow \beta A'$ et $A' \rightarrow \alpha A' \mid \alpha$
      \item \emph{Poussée vers l'avant.} Remplacer $A$ par sa définition dans les règles $B \rightarrow A \gamma$.
      \end{enumerate}
    \end{enumerate}
    
  \item On obtient la grammaire suivante :
    $$\left\langle \{a,b,c\}, \{A,B,C, A', B', C'\}, A,
    \left\{\begin{array}{rcl}
    A & \rightarrow & a A'\\
    B & \rightarrow & b B'\\
    C & \rightarrow & c C'\\
    A' & \rightarrow & b B'A' \mid b B' \\
    B' & \rightarrow & c C'B' \mid c C' \\
    C' & \rightarrow & a A'C' \mid a A' \\
    \end{array}\right\}\right\rangle$$

  \end{question}

\end{correction}

\endgroup
\endinput
